% Ad Familiares 9.14 (ad-familiares-09-14)
% Source: https://raw.githubusercontent.com/PerseusDL/canonical-latinLit/master/data/phi0474/phi056/phi0474.phi056.perseus-lat1.xml
% Fetched by scripts/fetch_latin.py

% Dateline (Perseus): Scr. in Pompeiano v Non. Mai. a. 710 (44).

\ciceroLetterOpener{CICERO DOLABELLAE CONSVLI SVO S.}

\ciceroSection{1} etsi contentus eram, mi Dolabella, tua gloria satisque ex ea magnam laetitiam voluptatemque capiebam, tamen non possum non confiteri cumulari me maximo gaudio, quod vulgo hominum opinio socium me adscribat tuis laudibus. neminem conveni (convenio autem cotidie plurimos ; sunt enim permulti optimi viri, qui valetudinis causa in haec loca veniant, praeterea ex municipiis frequentes necessarii mei), quin omnes, cum te summis laudibus ad caelum extulerunt, mihi continuo maximas gratias agant ; negant enim se dubitare quin tu meis praeceptis et consiliis obtemperans praestantissimum te civem et singularem consulem praebeas.

\ciceroSection{2} quibus ego quamquam verissime possum respondere te, quae facias, tuo iudicio et tua sponte facere nec cuiusquam egere consilio, tamen neque plane adsentior, ne imminuam tuam laudem, si omnis a meis consiliis profecta videatur, neque valde nego ; sum enim avidior etiam quam satis est gloriae ; et tamen non alienum est dignitate tua, quod ipsi Agamemnoni, regum regi, fuit honestum, habere aliquem in consiliis capiendis Nestorem, mihi vero gloriosum te iuvenem consulem florere laudibus quasi alumnum disciplinae meae.

\ciceroSection{3} L. quidem Caesar, cum ad eum aegrotum Neapolim venissem, quamquam erat oppressus totius corporis doloribus, tamen, ante quam me plane salutavit, 'O mi Cicero,' inquit, 'gratulor tibi, cum tantum vales apud Dolabellam, quantum si ego apud sororis filium valerem, iam salvi esse possemus ; Dolabellae vero tuo et gratulor et gratias ago, quem quidem post te consulem solum possumus vere consulem dicere.' deinde multa de facto ac de re gesta tua: nihil magnificentius, nihil praeclarius actum umquam, nihil rei p. salutarius. atque haec una vox omnium est.

\ciceroSection{4} A te autem peto ut me hanc quasi falsam hereditatem alienae gloriae sinas cernere meque aliqua ex parte in societatem tuarum laudum venire patiare. quamquam , mi Dolabella (haec enim iocatus sum) libentius omnis meas, si modo sunt aliquae meae laudes, ad te transfuderim quam aliquam partem exhauserim ex tuis. nam cum te semper tantum dilexerim, quantum tu intellegere potuisti, tum his tuis factis sic incensus sum, ut nihil umquam in amore fuerit ardentius. nihil est enim, mihi crede, virtute formosius, nihil pulchrius, nihil amabilius.

\ciceroSection{5} semper amavi, ut scis, M. Brutum propter eius summum ingenium, suavissimos mores, singularem probitatem atque constantiam ; tamen Idibus Martiis tantum accessit ad amorem ut mirarer locum fuisse augendi in eo, quod mihi iam pridem cumulatum etiam videbatur. quis erat qui putaret ad eum amorem, quem erga te habebam, pesse aliquid accedere? tantum accessit ut mihi nunc denique amare videar, antea dilexisse.

\ciceroSection{6} qua re quid est quod ego te horter ut dignitati et gloriae servias? proponam tibi claros viros, quod facere solent qui hortantur? neminem habeo clariorem quam te ipsum ; te imitere oportet, tecum ipse certes ; ne licet quidem tibi iam tantis rebus gestis non tui similem esse.

\ciceroSection{7} quod cum ita sit, hortatio non est necessaria, gratulatione magis utendum est ; contigit enim tibi quod haud scio an nemini, ut' summa severitas animadversionis non modo non invidiosa sed etiam popularis esset et cum bonis omnibus tum infimo cuique gratissima. hoc si tibi fortuna quadam contigisset, gratularer felicitati tuae ; sed contigit magnitudine cum animi tum etiam ingeni atque consili. legi enim contionem tuam ; nihil illa sapientius ; ita pedetemptim et gradatim tum accessus a te ad causam facti, tum recessus, ut res ipsa maturitatem tibi animadvertendi omnium concessu daret.

\ciceroSection{8} liberasti igitur et urbem periculo et civitatem metu neque solum ad tempus maximam utilitatem attulisti sed etiam ad exemplum. quo facto intellegere debes in te positam esse rem p. tibique non modo tuendos, sed etiam ornandos esse illos viros, a quibus initium libertatis profectum est. sed his de rebus coram plura propediem, ut spero. tu quoniam rem p. nosque conservas, fac ut diligentissime te ipsum, mi Dolabella, custodias.
